\Chapter{98}

\Verse{1}
{
    \zR{话说宝玉见了贾政，回至房中，更觉头昏脑闷，懒怠动弹，连饭也没吃，便昏 沉睡去。}
}
{
    \eR{[English source not present in the checked-in Bencraft/Joly files.]}
}
{
}

\Verse{2}
{
    \zR{仍旧延医诊治，服药不效，索性连人也认不明白了。}
}
{
    \eR{[English source not present in the checked-in Bencraft/Joly files.]}
}
{
}

\Verse{3}
{
    \zR{大家扶着他坐起来， 还是像个好人。}
}
{
    \eR{[English source not present in the checked-in Bencraft/Joly files.]}
}
{
}

\Verse{4}
{
    \zR{一连闹了几天。}
}
{
    \eR{[English source not present in the checked-in Bencraft/Joly files.]}
}
{
}

\Verse{5}
{
    \zR{那日恰是回九之期，说是若不过去，薛姨妈脸上过 不去；若说去呢，宝玉这般光景，明知是为黛玉而起，欲要告诉明白，又恐气急生 变。}
}
{
    \eR{[English source not present in the checked-in Bencraft/Joly files.]}
}
{
}

\Verse{6}
{
    \zR{宝钗是新媳妇，又难劝慰，必得姨妈过来才好。}
}
{
    \eR{[English source not present in the checked-in Bencraft/Joly files.]}
}
{
}

\Verse{7}
{
    \zR{若不回九，姨妈嗔怪。}
}
{
    \eR{[English source not present in the checked-in Bencraft/Joly files.]}
}
{
}

\Verse{8}
{
    \zR{便与王 夫人凤姐商议道：“我看宝玉竟是魂不守舍，起动是不怕的。}
}
{
    \eR{[English source not present in the checked-in Bencraft/Joly files.]}
}
{
}

\Verse{9}
{
    \zR{用两乘小轿，叫人扶 着，从园里过去，应了回九的吉期；以后请姨妈过来安慰宝钗，咱们一心一计的调 治宝玉，可不两全？”}
}
{
    \eR{[English source not present in the checked-in Bencraft/Joly files.]}
}
{
}

\Verse{10}
{
    \zR{王夫人答应了，即刻预备。}
}
{
    \eR{[English source not present in the checked-in Bencraft/Joly files.]}
}
{
}

\Verse{11}
{
    \zR{幸亏宝钗是新媳妇，宝玉是个疯 傻的，由人掇弄过去了，宝钗也明知其事，心里只怨母亲办得糊涂，事已至此，不 肯多言。}
}
{
    \eR{[English source not present in the checked-in Bencraft/Joly files.]}
}
{
}

\Verse{12}
{
    \zR{独有薛姨妈看见宝玉这般光景，心里懊悔，只得草草完事。}
}
{
    \eR{[English source not present in the checked-in Bencraft/Joly files.]}
}
{
}

\Verse{13}
{
    \zR{回家，宝玉越加沉重。}
}
{
    \eR{[English source not present in the checked-in Bencraft/Joly files.]}
}
{
}

\Verse{14}
{
    \zR{次日连起坐都不能了，日重一日，甚至汤水不进。}
}
{
    \eR{[English source not present in the checked-in Bencraft/Joly files.]}
}
{
}

\Verse{15}
{
    \zR{薛姨 妈等忙了手脚，各处遍请名医，皆不识病源。}
}
{
    \eR{[English source not present in the checked-in Bencraft/Joly files.]}
}
{
}

\Verse{16}
{
    \zR{只有城外破寺中住着个穷医姓毕别号 知庵的，诊得病源是悲喜激射，冷暖失调，饮食失时，忧忿滞中，正气壅闭：此内 伤外感之症。}
}
{
    \eR{[English source not present in the checked-in Bencraft/Joly files.]}
}
{
}

\Verse{17}
{
    \zR{于是度量用药。}
}
{
    \eR{[English source not present in the checked-in Bencraft/Joly files.]}
}
{
}

\Verse{18}
{
    \zR{至晚服了，二更后，果然省些人事，便要喝水。}
}
{
    \eR{[English source not present in the checked-in Bencraft/Joly files.]}
}
{
}

\Verse{19}
{
    \zR{贾母 王夫人等才放了心，请了薛姨妈带了宝钗，都到贾母那里，暂且歇息。}
}
{
    \eR{[English source not present in the checked-in Bencraft/Joly files.]}
}
{
}

\Verse{20}
{
    \zR{宝玉片时清 楚，自料难保，见诸人散后，房中只有袭人，因唤袭人至跟前，拉着手哭道：“我 问你：宝姐姐怎么来的？}
}
{
    \eR{[English source not present in the checked-in Bencraft/Joly files.]}
}
{
}

\Verse{21}
{
    \zR{我记得老爷给我娶了林妹妹过来，怎么叫宝姐姐赶出去了？}
}
{
    \eR{[English source not present in the checked-in Bencraft/Joly files.]}
}
{
}

\Verse{22}
{
    \zR{他为什么霸占住在这里？}
}
{
    \eR{[English source not present in the checked-in Bencraft/Joly files.]}
}
{
}

\Verse{23}
{
    \zR{我要说呢，又恐怕得罪了他。}
}
{
    \eR{[English source not present in the checked-in Bencraft/Joly files.]}
}
{
}

\Verse{24}
{
    \zR{你们听见林妹妹哭的怎么样 了？”}
}
{
    \eR{[English source not present in the checked-in Bencraft/Joly files.]}
}
{
}

\Verse{25}
{
    \zR{袭人不敢明说，只得说道：“林姑娘病着呢。”}
}
{
    \eR{[English source not present in the checked-in Bencraft/Joly files.]}
}
{
}

\Verse{26}
{
    \zR{宝玉又道：“我瞧瞧他去。”}
}
{
    \eR{[English source not present in the checked-in Bencraft/Joly files.]}
}
{
}

\Verse{27}
{
    \zR{说着要起来。}
}
{
    \eR{[English source not present in the checked-in Bencraft/Joly files.]}
}
{
}

\Verse{28}
{
    \zR{那知连日饮食不进，身子岂能动转？}
}
{
    \eR{[English source not present in the checked-in Bencraft/Joly files.]}
}
{
}

\Verse{29}
{
    \zR{便哭道：“我要死了!我有一句心 里的话，只求你回明老太太：横竖林妹妹也是要死的，我如今也不能保，两处两个 病人，都要死的。}
}
{
    \eR{[English source not present in the checked-in Bencraft/Joly files.]}
}
{
}

\Verse{30}
{
    \zR{死了越发难张罗，不如腾一处空房子，趁早把我和林妹妹两个抬 在那里，活着也好一处医治、伏侍，死了也好一处停放。}
}
{
    \eR{[English source not present in the checked-in Bencraft/Joly files.]}
}
{
}

\Verse{31}
{
    \zR{你依我这话，不枉了几年 的情分。”}
}
{
    \eR{[English source not present in the checked-in Bencraft/Joly files.]}
}
{
}

\Verse{32}
{
    \zR{袭人听了这些话，又急，又笑，又痛。}
}
{
    \eR{[English source not present in the checked-in Bencraft/Joly files.]}
}
{
}

\Verse{33}
{
    \zR{宝钗恰好同着莺儿过来，也听见了。}
}
{
    \eR{[English source not present in the checked-in Bencraft/Joly files.]}
}
{
}

\Verse{34}
{
    \zR{便说道：“你放着病不保养，何苦说这些 不吉利的话呢？}
}
{
    \eR{[English source not present in the checked-in Bencraft/Joly files.]}
}
{
}

\Verse{35}
{
    \zR{老太太才安慰了些，你又生出事来。}
}
{
    \eR{[English source not present in the checked-in Bencraft/Joly files.]}
}
{
}

\Verse{36}
{
    \zR{老太太一生疼你一个，如今八 十多岁的人了，虽不图你的诰封，将来你成了人，老太太也看着乐一天，也不枉了 老人家的苦心。}
}
{
    \eR{[English source not present in the checked-in Bencraft/Joly files.]}
}
{
}

\Verse{37}
{
    \zR{太太更是不必说了，一生的心血精神，抚养了你这一个儿子，若是 半途死了，太太将来怎么样呢？}
}
{
    \eR{[English source not present in the checked-in Bencraft/Joly files.]}
}
{
}

\Verse{38}
{
    \zR{我虽是薄命，也不至于此。}
}
{
    \eR{[English source not present in the checked-in Bencraft/Joly files.]}
}
{
}

\Verse{39}
{
    \zR{据此三件看来，你就要 死，那天也不容你死的，所以你是不能死的。}
}
{
    \eR{[English source not present in the checked-in Bencraft/Joly files.]}
}
{
}

\Verse{40}
{
    \zR{只管安稳着养个四五天后，风邪散了， 太和正气一足，自然这些邪病都没有了。”}
}
{
    \eR{[English source not present in the checked-in Bencraft/Joly files.]}
}
{
}

\Verse{41}
{
    \zR{宝玉听了，竟是无言可答，半晌，方才 嘻嘻的笑道：“你是好些时不和我说话了，这会子说这些大道理的话给谁听？”}
}
{
    \eR{[English source not present in the checked-in Bencraft/Joly files.]}
}
{
}

\Verse{42}
{
    \zR{宝 钗听了这话，便又说道：“实告诉你说罢：那两日你不知人事的时候，林妹妹已经 亡故了！”}
}
{
    \eR{[English source not present in the checked-in Bencraft/Joly files.]}
}
{
}

\Verse{43}
{
    \zR{宝玉忽然坐起，大声诧异道：“果真死了吗？”}
}
{
    \eR{[English source not present in the checked-in Bencraft/Joly files.]}
}
{
}

\Verse{44}
{
    \zR{宝钗道：“果真死了， 岂有红口白舌咒人死的呢!老太太、太太知道你姐妹和睦，你听见他死了，自然你 也要死，所以不肯告诉你。”}
}
{
    \eR{[English source not present in the checked-in Bencraft/Joly files.]}
}
{
}

\Verse{45}
{
    \zR{宝玉听了，不禁放声大哭，倒在床上，忽然眼前漆黑，辨不出方向。}
}
{
    \eR{[English source not present in the checked-in Bencraft/Joly files.]}
}
{
}

\Verse{46}
{
    \zR{心中正自 恍惚，只见眼前好像有人走来。}
}
{
    \eR{[English source not present in the checked-in Bencraft/Joly files.]}
}
{
}

\Verse{47}
{
    \zR{宝玉茫然问道：“借问此是何处？”}
}
{
    \eR{[English source not present in the checked-in Bencraft/Joly files.]}
}
{
}

\Verse{48}
{
    \zR{那人道：“此 阴司泉路。}
}
{
    \eR{[English source not present in the checked-in Bencraft/Joly files.]}
}
{
}

\Verse{49}
{
    \zR{你寿未终，何故至此？”}
}
{
    \eR{[English source not present in the checked-in Bencraft/Joly files.]}
}
{
}

\Verse{50}
{
    \zR{宝玉道：“适闻有一故人已死，遂寻访至此， 不觉迷途。”}
}
{
    \eR{[English source not present in the checked-in Bencraft/Joly files.]}
}
{
}

\Verse{51}
{
    \zR{那人道：“故人是谁？”}
}
{
    \eR{[English source not present in the checked-in Bencraft/Joly files.]}
}
{
}

\Verse{52}
{
    \zR{宝玉道：“姑苏林黛玉。”}
}
{
    \eR{[English source not present in the checked-in Bencraft/Joly files.]}
}
{
}

\Verse{53}
{
    \zR{那人冷笑道：“林 黛玉生不同人，死不同鬼，无魂无魄，何处寻访？}
}
{
    \eR{[English source not present in the checked-in Bencraft/Joly files.]}
}
{
}

\Verse{54}
{
    \zR{凡人魂魄，聚而成形，散而为气， 生前聚之，死则散焉。}
}
{
    \eR{[English source not present in the checked-in Bencraft/Joly files.]}
}
{
}

\Verse{55}
{
    \zR{常人尚无可寻访，何况林黛玉呢？}
}
{
    \eR{[English source not present in the checked-in Bencraft/Joly files.]}
}
{
}

\Verse{56}
{
    \zR{汝快回去罢。”}
}
{
    \eR{[English source not present in the checked-in Bencraft/Joly files.]}
}
{
}

\Verse{57}
{
    \zR{宝玉听了， 呆了半晌，道：“既云死者散也，又如何有这个阴司呢？”}
}
{
    \eR{[English source not present in the checked-in Bencraft/Joly files.]}
}
{
}

\Verse{58}
{
    \zR{那人冷笑道：“那阴司， 说有便有，说无就无。}
}
{
    \eR{[English source not present in the checked-in Bencraft/Joly files.]}
}
{
}

\Verse{59}
{
    \zR{皆为世俗溺于生死之说，设言以警世，便道上天深怒愚人： 或不守分安常；或生禄未终，自行夭折；或嗜淫欲，尚气逞凶，无故自殒者，特设
        此地狱，囚其魂魄，受无边的苦，以偿生前之罪，汝寻黛玉，是无故自陷也。}
}
{
    \eR{[English source not present in the checked-in Bencraft/Joly files.]}
}
{
}

\Verse{60}
{
    \zR{且黛 玉已归太虚幻境，汝若有心寻访，潜心修养，自然有时相见；如不安生，即以自行 夭折之罪，囚禁阴司，除父母之外，图一见黛玉，终不能矣。”}
}
{
    \eR{[English source not present in the checked-in Bencraft/Joly files.]}
}
{
}

\Verse{61}
{
    \zR{那人说毕，袖中取 出一石，向宝玉心口掷来。}
}
{
    \eR{[English source not present in the checked-in Bencraft/Joly files.]}
}
{
}

\Verse{62}
{
    \zR{宝玉听了这话，又被这石子打着心窝，吓的即欲回家， 只恨迷了道路。}
}
{
    \eR{[English source not present in the checked-in Bencraft/Joly files.]}
}
{
}

\Verse{63}
{
    \zR{正在踌躇，忽听那边有人唤他。}
}
{
    \eR{[English source not present in the checked-in Bencraft/Joly files.]}
}
{
}

\Verse{64}
{
    \zR{回首看时，不是别人，正是贾母、 王夫人、宝钗、袭人等围绕哭泣叫着，自己仍旧躺在床上。}
}
{
    \eR{[English source not present in the checked-in Bencraft/Joly files.]}
}
{
}

\Verse{65}
{
    \zR{见案上红灯，窗前皓月， 依然锦绣丛中，繁华世界。}
}
{
    \eR{[English source not present in the checked-in Bencraft/Joly files.]}
}
{
}

\Verse{66}
{
    \zR{定神一想，原来竟是一场大梦。}
}
{
    \eR{[English source not present in the checked-in Bencraft/Joly files.]}
}
{
}

\Verse{67}
{
    \zR{浑身冷汗，觉得心内清 爽。}
}
{
    \eR{[English source not present in the checked-in Bencraft/Joly files.]}
}
{
}

\Verse{68}
{
    \zR{仔细一想，真正无可奈何，不过长叹数声。}
}
{
    \eR{[English source not present in the checked-in Bencraft/Joly files.]}
}
{
}

\Verse{69}
{
    \zR{起初宝钗早知黛玉已死，因贾母等不许众人告诉宝玉知道，恐添病难治。}
}
{
    \eR{[English source not present in the checked-in Bencraft/Joly files.]}
}
{
}

\Verse{70}
{
    \zR{自己 却深知宝玉之病实因黛玉而起，失玉次之，故趁势说明，使其一痛决绝，神魂一归， 庶可疗治。}
}
{
    \eR{[English source not present in the checked-in Bencraft/Joly files.]}
}
{
}

\Verse{71}
{
    \zR{贾母王夫人等不知宝钗的用意，深怪他造次，后来见宝玉醒了过来，方 才放心，立刻到外书房请了毕大夫进来诊视。}
}
{
    \eR{[English source not present in the checked-in Bencraft/Joly files.]}
}
{
}

\Verse{72}
{
    \zR{那大夫进来诊了脉，便道奇怪：“这 回脉气沉静，神安郁散，明日进调理的药，就可以望好了。”}
}
{
    \eR{[English source not present in the checked-in Bencraft/Joly files.]}
}
{
}

\Verse{73}
{
    \zR{说着出去。}
}
{
    \eR{[English source not present in the checked-in Bencraft/Joly files.]}
}
{
}

\Verse{74}
{
    \zR{众人各自 安心散去。}
}
{
    \eR{[English source not present in the checked-in Bencraft/Joly files.]}
}
{
}

\Verse{75}
{
    \zR{袭人起初深怨宝钗不该告诉，惟是口中不好说出。}
}
{
    \eR{[English source not present in the checked-in Bencraft/Joly files.]}
}
{
}

\Verse{76}
{
    \zR{莺儿背地也说宝钗道： “姑娘忒性急了。”}
}
{
    \eR{[English source not present in the checked-in Bencraft/Joly files.]}
}
{
}

\Verse{77}
{
    \zR{宝钗道：“你知道什么!好歹横竖有我呢。”}
}
{
    \eR{[English source not present in the checked-in Bencraft/Joly files.]}
}
{
}

\Verse{78}
{
    \zR{那宝钗任人诽谤，并不介意，只窥察宝玉心病，暗下针砭。}
}
{
    \eR{[English source not present in the checked-in Bencraft/Joly files.]}
}
{
}

\Verse{79}
{
    \zR{一日，宝玉渐觉神 志安定，虽一时想起黛玉，尚有糊涂。}
}
{
    \eR{[English source not present in the checked-in Bencraft/Joly files.]}
}
{
}

\Verse{80}
{
    \zR{更有袭人缓缓的将“老爷选定的宝姑娘为人 和厚，嫌林姑娘秉性古怪，原恐早夭。}
}
{
    \eR{[English source not present in the checked-in Bencraft/Joly files.]}
}
{
}

\Verse{81}
{
    \zR{老太太恐你不知好歹，病中着急，所以叫雪 雁过来哄你”的话，时常劝解。}
}
{
    \eR{[English source not present in the checked-in Bencraft/Joly files.]}
}
{
}

\Verse{82}
{
    \zR{宝玉终是心酸落泪。}
}
{
    \eR{[English source not present in the checked-in Bencraft/Joly files.]}
}
{
}

\Verse{83}
{
    \zR{欲待寻死，又想着梦中之言， 又恐老太太、太太生气，又不得撩开。}
}
{
    \eR{[English source not present in the checked-in Bencraft/Joly files.]}
}
{
}

\Verse{84}
{
    \zR{又想黛玉已死，宝钗又是第一等人物，方信 “金石姻缘”有定，自己也解了好些。}
}
{
    \eR{[English source not present in the checked-in Bencraft/Joly files.]}
}
{
}

\Verse{85}
{
    \zR{宝钗看来不妨大事，于是自己心也安了，只 在贾母王夫人等前尽行过家庭之礼后，便设法以释宝玉之忧。}
}
{
    \eR{[English source not present in the checked-in Bencraft/Joly files.]}
}
{
}

\Verse{86}
{
    \zR{宝玉虽不能时常坐 起，亦常见宝钗坐在床前，禁不住生来旧病。}
}
{
    \eR{[English source not present in the checked-in Bencraft/Joly files.]}
}
{
}

\Verse{87}
{
    \zR{宝钗每以正言解劝，以“养身要紧， 你我既为夫妇，岂在一时”之语安慰他。}
}
{
    \eR{[English source not present in the checked-in Bencraft/Joly files.]}
}
{
}

\Verse{88}
{
    \zR{那宝玉心里虽不顺遂，无奈日里贾母王夫 人及薛姨妈等轮流相伴，夜间宝钗独去安寝，贾母又派人服侍，只得安心静养。}
}
{
    \eR{[English source not present in the checked-in Bencraft/Joly files.]}
}
{
}

\Verse{89}
{
    \zR{又 见宝钗举动温柔，就也渐渐的将爱慕黛玉的心肠略移在宝钗身上。}
}
{
    \eR{[English source not present in the checked-in Bencraft/Joly files.]}
}
{
}

\Verse{90}
{
    \zR{此是后话。}
}
{
    \eR{[English source not present in the checked-in Bencraft/Joly files.]}
}
{
}

\Verse{91}
{
    \zR{却说宝玉成家的那一日，黛玉白日已经昏晕过去，却心头口中一丝微气不断， 把个李纨和紫鹃哭的死去活来。}
}
{
    \eR{[English source not present in the checked-in Bencraft/Joly files.]}
}
{
}

\Verse{92}
{
    \zR{到了晚间，黛玉却又缓过来了，微微睁开眼，似有 要水要汤的光景。}
}
{
    \eR{[English source not present in the checked-in Bencraft/Joly files.]}
}
{
}

\Verse{93}
{
    \zR{此时雪雁已去，只有紫鹃和李纨在旁。}
}
{
    \eR{[English source not present in the checked-in Bencraft/Joly files.]}
}
{
}

\Verse{94}
{
    \zR{紫鹃便端了一盏桂圆汤和 的梨汁，用小银匙灌了两三匙。}
}
{
    \eR{[English source not present in the checked-in Bencraft/Joly files.]}
}
{
}

\Verse{95}
{
    \zR{黛玉闭着眼，静养了一会子，觉得心里似明似暗的。}
}
{
    \eR{[English source not present in the checked-in Bencraft/Joly files.]}
}
{
}

\Verse{96}
{
    \zR{此时李纨见黛玉略缓，明知是回光返照的光景，却料着还有一半天耐头，自己回到 稻香村，料理了一回事情。}
}
{
    \eR{[English source not present in the checked-in Bencraft/Joly files.]}
}
{
}

\Verse{97}
{
    \zR{这里黛玉睁开眼一看，只有紫鹃和奶妈并几个小丫头在那里，便一手攥了紫鹃 的手，使着劲说道：“我是不中用的人了!你伏侍我几年，我原指望咱们两个总在
        一处，不想我――”说着，又喘了一会儿，闭了眼歇着。}
}
{
    \eR{[English source not present in the checked-in Bencraft/Joly files.]}
}
{
}

\Verse{98}
{
    \zR{紫鹃见他攥着不肯松手， 自己也不敢挪动。}
}
{
    \eR{[English source not present in the checked-in Bencraft/Joly files.]}
}
{
}

\Verse{99}
{
    \zR{看他的光景，比早半天好些，只当还可以回转，听了这话，又寒 了半截。}
}
{
    \eR{[English source not present in the checked-in Bencraft/Joly files.]}
}
{
}

\Verse{100}
{
    \zR{半天，黛玉又说道：“妹妹!我这里并没亲人，我的身子是干净的，你好 歹叫他们送我回去。”}
}
{
    \eR{[English source not present in the checked-in Bencraft/Joly files.]}
}
{
}

\Verse{101}
{
    \zR{说到这里，又闭了眼不言语了。}
}
{
    \eR{[English source not present in the checked-in Bencraft/Joly files.]}
}
{
}

\Verse{102}
{
    \zR{那手却渐渐紧了，喘成一处， 只是出气大，入气小，已经促疾的很了。}
}
{
    \eR{[English source not present in the checked-in Bencraft/Joly files.]}
}
{
}

\Verse{103}
{
    \zR{紫鹃忙了，连忙叫人请李纨。}
}
{
    \eR{[English source not present in the checked-in Bencraft/Joly files.]}
}
{
}

\Verse{104}
{
    \zR{可巧探春来了。}
}
{
    \eR{[English source not present in the checked-in Bencraft/Joly files.]}
}
{
}

\Verse{105}
{
    \zR{紫鹃见了，忙悄悄的说道：“三 姑娘，瞧瞧林姑娘罢。”}
}
{
    \eR{[English source not present in the checked-in Bencraft/Joly files.]}
}
{
}

\Verse{106}
{
    \zR{说着，泪如雨下。}
}
{
    \eR{[English source not present in the checked-in Bencraft/Joly files.]}
}
{
}

\Verse{107}
{
    \zR{探春过来，摸了摸黛玉的手，已经凉了， 连目光也都散了。}
}
{
    \eR{[English source not present in the checked-in Bencraft/Joly files.]}
}
{
}

\Verse{108}
{
    \zR{探春紫鹃正哭着叫人端水来给黛玉擦洗，李纨赶忙进来了。}
}
{
    \eR{[English source not present in the checked-in Bencraft/Joly files.]}
}
{
}

\Verse{109}
{
    \zR{三个 人才见了，不及说话。}
}
{
    \eR{[English source not present in the checked-in Bencraft/Joly files.]}
}
{
}

\Verse{110}
{
    \zR{刚擦着，猛听黛玉直声叫道：“宝玉!宝玉!你好――”说到 “好”字，便浑身冷汗，不作声了。}
}
{
    \eR{[English source not present in the checked-in Bencraft/Joly files.]}
}
{
}

\Verse{111}
{
    \zR{紫鹃等急忙扶住，那汗愈出，身子便渐渐的冷 了。}
}
{
    \eR{[English source not present in the checked-in Bencraft/Joly files.]}
}
{
}

\Verse{112}
{
    \zR{探春李纨叫人乱着拢头穿衣，只见黛玉两眼一翻，呜呼! 香魂一缕随风散，愁绪三更入梦遥! 当时黛玉气绝，正是宝玉娶宝钗的这个时辰。}
}
{
    \eR{[English source not present in the checked-in Bencraft/Joly files.]}
}
{
}

\Verse{113}
{
    \zR{紫鹃等都大哭起来。}
}
{
    \eR{[English source not present in the checked-in Bencraft/Joly files.]}
}
{
}

\Verse{114}
{
    \zR{李纨探春想他素 日的可疼，今日更加可怜，便也伤心痛哭。}
}
{
    \eR{[English source not present in the checked-in Bencraft/Joly files.]}
}
{
}

\Verse{115}
{
    \zR{因潇湘馆离新房子甚远，所以那边并没 听见。}
}
{
    \eR{[English source not present in the checked-in Bencraft/Joly files.]}
}
{
}

\Verse{116}
{
    \zR{一时，大家痛哭了一阵，只听得远远一阵音乐之声，侧耳一听，却又没有了。}
}
{
    \eR{[English source not present in the checked-in Bencraft/Joly files.]}
}
{
}

\Verse{117}
{
    \zR{探春李纨走出院外再听时，惟有竹梢风动，月影移墙，好不凄凉冷淡。}
}
{
    \eR{[English source not present in the checked-in Bencraft/Joly files.]}
}
{
}

\Verse{118}
{
    \zR{一时叫了林之孝家的过来，将黛玉停放毕，派人看守，等明早去回凤姐。}
}
{
    \eR{[English source not present in the checked-in Bencraft/Joly files.]}
}
{
}

\Verse{119}
{
    \zR{凤姐 因见贾母王夫人等忙乱，贾政起身，又为宝玉昏愦更甚，正在着急异常之时，若是 又将黛玉的凶信回了，恐贾母王夫人愁苦交加，急出病来，只得亲自到园。}
}
{
    \eR{[English source not present in the checked-in Bencraft/Joly files.]}
}
{
}

\Verse{120}
{
    \zR{到了潇 湘馆内，也不免哭了一场。}
}
{
    \eR{[English source not present in the checked-in Bencraft/Joly files.]}
}
{
}

\Verse{121}
{
    \zR{见了李纨探春，知道诸事齐备，就说：“很好。}
}
{
    \eR{[English source not present in the checked-in Bencraft/Joly files.]}
}
{
}

\Verse{122}
{
    \zR{只是刚 才你们为什么不言语，叫我着急？”}
}
{
    \eR{[English source not present in the checked-in Bencraft/Joly files.]}
}
{
}

\Verse{123}
{
    \zR{探春道：“刚才送老爷，怎么说呢？”}
}
{
    \eR{[English source not present in the checked-in Bencraft/Joly files.]}
}
{
}

\Verse{124}
{
    \zR{凤姐道： “这倒是你们两个可怜他些。}
}
{
    \eR{[English source not present in the checked-in Bencraft/Joly files.]}
}
{
}

\Verse{125}
{
    \zR{这么着，我还得那边去招呼那个冤家呢。}
}
{
    \eR{[English source not present in the checked-in Bencraft/Joly files.]}
}
{
}

\Verse{126}
{
    \zR{但是这件事 好累坠：若是今日不回，使不得；若回了，恐怕老太太搁不住。”}
}
{
    \eR{[English source not present in the checked-in Bencraft/Joly files.]}
}
{
}

\Verse{127}
{
    \zR{李纨道：“你去 见机行事，得回再回方好。”}
}
{
    \eR{[English source not present in the checked-in Bencraft/Joly files.]}
}
{
}

\Verse{128}
{
    \zR{凤姐点头，忙忙的去了。}
}
{
    \eR{[English source not present in the checked-in Bencraft/Joly files.]}
}
{
}

\Verse{129}
{
    \zR{凤姐到了宝玉那里，听见大夫说不妨事，贾母王夫人略觉放心，凤姐便背了宝 玉，缓缓的将黛玉的事回明了。}
}
{
    \eR{[English source not present in the checked-in Bencraft/Joly files.]}
}
{
}

\Verse{130}
{
    \zR{贾母王夫人听得，都唬了一大跳。}
}
{
    \eR{[English source not present in the checked-in Bencraft/Joly files.]}
}
{
}

\Verse{131}
{
    \zR{贾母眼泪交流， 说道：“是我弄坏了他了。}
}
{
    \eR{[English source not present in the checked-in Bencraft/Joly files.]}
}
{
}

\Verse{132}
{
    \zR{但只是这个丫头也忒傻气！”}
}
{
    \eR{[English source not present in the checked-in Bencraft/Joly files.]}
}
{
}

\Verse{133}
{
    \zR{说着，便要到园里去哭他 一场，又惦记着宝玉，两头难顾。}
}
{
    \eR{[English source not present in the checked-in Bencraft/Joly files.]}
}
{
}

\Verse{134}
{
    \zR{王夫人等含悲共劝贾母：“不必过去，老太太身 子要紧。”}
}
{
    \eR{[English source not present in the checked-in Bencraft/Joly files.]}
}
{
}

\Verse{135}
{
    \zR{贾母无奈，只得叫王夫人自去。}
}
{
    \eR{[English source not present in the checked-in Bencraft/Joly files.]}
}
{
}

\Verse{136}
{
    \zR{又说：“你替我告诉他的阴灵：‘并不 是我忍心不来送你，只为有个亲疏。}
}
{
    \eR{[English source not present in the checked-in Bencraft/Joly files.]}
}
{
}

\Verse{137}
{
    \zR{你是我的外孙女儿，是亲的了；若与宝玉比起 来，可是宝玉比你更亲些。}
}
{
    \eR{[English source not present in the checked-in Bencraft/Joly files.]}
}
{
}

\Verse{138}
{
    \zR{倘宝玉有些不好，我怎么见他父亲呢！’”}
}
{
    \eR{[English source not present in the checked-in Bencraft/Joly files.]}
}
{
}

\Verse{139}
{
    \zR{说着，又哭 起来。}
}
{
    \eR{[English source not present in the checked-in Bencraft/Joly files.]}
}
{
}

\Verse{140}
{
    \zR{王夫人劝道：“林姑娘是老太太最疼的，但只寿夭有定，如今已经死了，无 可尽心，只是葬礼上要上等的发送。}
}
{
    \eR{[English source not present in the checked-in Bencraft/Joly files.]}
}
{
}

\Verse{141}
{
    \zR{一则可以少尽咱们的心，二则就是姑太太和外 甥女儿的阴灵儿也可以少安了。”}
}
{
    \eR{[English source not present in the checked-in Bencraft/Joly files.]}
}
{
}

\Verse{142}
{
    \zR{贾母听到这里，越发痛哭起来。}
}
{
    \eR{[English source not present in the checked-in Bencraft/Joly files.]}
}
{
}

\Verse{143}
{
    \zR{凤姐恐怕老人家 伤感太过，明仗着宝玉心中不甚明白，便偷偷的使人来撒个谎儿，哄老太太道：“宝 玉那里找老太太呢。”}
}
{
    \eR{[English source not present in the checked-in Bencraft/Joly files.]}
}
{
}

\Verse{144}
{
    \zR{贾母听见，才止住泪问道：“不是又有什么缘故？”}
}
{
    \eR{[English source not present in the checked-in Bencraft/Joly files.]}
}
{
}

\Verse{145}
{
    \zR{凤姐陪 笑道：“没什么缘故，他大约是想老太太的意思。”}
}
{
    \eR{[English source not present in the checked-in Bencraft/Joly files.]}
}
{
}

\Verse{146}
{
    \zR{贾母连忙扶了珍珠儿，凤姐也 跟着过来。}
}
{
    \eR{[English source not present in the checked-in Bencraft/Joly files.]}
}
{
}

\Verse{147}
{
    \zR{走至半路，正遇王夫人过来，一一回明了贾母，贾母自然又是哀痛的； 只因要到宝玉那边，只得忍泪含悲的说道：“既这么着，我也不过去了，由你们办 罢。}
}
{
    \eR{[English source not present in the checked-in Bencraft/Joly files.]}
}
{
}

\Verse{148}
{
    \zR{我看着心里也难受，只别委屈了他就是了。”}
}
{
    \eR{[English source not present in the checked-in Bencraft/Joly files.]}
}
{
}

\Verse{149}
{
    \zR{王夫人凤姐一一答应了，贾母才 过宝玉这边来。}
}
{
    \eR{[English source not present in the checked-in Bencraft/Joly files.]}
}
{
}

\Verse{150}
{
    \zR{见了宝玉，因问：“你做什么找我？”}
}
{
    \eR{[English source not present in the checked-in Bencraft/Joly files.]}
}
{
}

\Verse{151}
{
    \zR{宝玉笑道：“我昨日晚上看 见林妹妹来了，他说要回南去，我想没人留的住，还得老太太给我留一留他。”}
}
{
    \eR{[English source not present in the checked-in Bencraft/Joly files.]}
}
{
}

\Verse{152}
{
    \zR{贾 母听着，说：“使得，只管放心罢。”}
}
{
    \eR{[English source not present in the checked-in Bencraft/Joly files.]}
}
{
}

\Verse{153}
{
    \zR{袭人因扶宝玉躺下。}
}
{
    \eR{[English source not present in the checked-in Bencraft/Joly files.]}
}
{
}

\Verse{154}
{
    \zR{贾母出来，到宝钗这边 来。}
}
{
    \eR{[English source not present in the checked-in Bencraft/Joly files.]}
}
{
}

\Verse{155}
{
    \zR{那时宝钗尚未回九，所以每每见了人，倒有些含羞之意。}
}
{
    \eR{[English source not present in the checked-in Bencraft/Joly files.]}
}
{
}

\Verse{156}
{
    \zR{这一天，见贾母满面 泪痕，递了茶，贾母叫他坐下。}
}
{
    \eR{[English source not present in the checked-in Bencraft/Joly files.]}
}
{
}

\Verse{157}
{
    \zR{宝钗侧身陪着坐了，才问道：“听得林妹妹病了， 不知他可好些了？”}
}
{
    \eR{[English source not present in the checked-in Bencraft/Joly files.]}
}
{
}

\Verse{158}
{
    \zR{贾母听了这话，那眼泪止不住流下来，因说道：“我的儿!我 告诉你，你可别告诉宝玉。}
}
{
    \eR{[English source not present in the checked-in Bencraft/Joly files.]}
}
{
}

\Verse{159}
{
    \zR{都是因你林妹妹，才叫你受了多少委屈!你如今作媳妇 了，我才告诉你：这如今你林妹妹没了两三天了，就是娶你的那个时辰死的。}
}
{
    \eR{[English source not present in the checked-in Bencraft/Joly files.]}
}
{
}

\Verse{160}
{
    \zR{如今 宝玉这一番病，还是为着这个。}
}
{
    \eR{[English source not present in the checked-in Bencraft/Joly files.]}
}
{
}

\Verse{161}
{
    \zR{你们先都在园子里，自然也都是明白的。”}
}
{
    \eR{[English source not present in the checked-in Bencraft/Joly files.]}
}
{
}

\Verse{162}
{
    \zR{宝钗把 脸飞红了，想到黛玉之死，又不免落下泪来。}
}
{
    \eR{[English source not present in the checked-in Bencraft/Joly files.]}
}
{
}

\Verse{163}
{
    \zR{贾母又说了一回话去了。}
}
{
    \eR{[English source not present in the checked-in Bencraft/Joly files.]}
}
{
}

\Verse{164}
{
    \zR{自此，宝钗千回万转，想了一个主意，只不肯造次，所以过了回九，才想出这 个法子来。}
}
{
    \eR{[English source not present in the checked-in Bencraft/Joly files.]}
}
{
}

\Verse{165}
{
    \zR{如今果然好些，然后大家说话才不至似前留神。}
}
{
    \eR{[English source not present in the checked-in Bencraft/Joly files.]}
}
{
}

\Verse{166}
{
    \zR{独是宝玉虽然病势一天 好似一天，他的痴心总不能解，必要亲去哭他一场。}
}
{
    \eR{[English source not present in the checked-in Bencraft/Joly files.]}
}
{
}

\Verse{167}
{
    \zR{贾母等知他病未除根，不许他 胡思乱想，怎奈他郁闷难堪，病多反复，倒是大夫看出心病，索性叫他开散了再用 药调理，倒可好得快些。}
}
{
    \eR{[English source not present in the checked-in Bencraft/Joly files.]}
}
{
}

\Verse{168}
{
    \zR{宝玉听说，立刻要往潇湘馆来。}
}
{
    \eR{[English source not present in the checked-in Bencraft/Joly files.]}
}
{
}

\Verse{169}
{
    \zR{贾母等只得叫人抬了竹椅 子过来，扶宝玉坐上，贾母王夫人即便先行。}
}
{
    \eR{[English source not present in the checked-in Bencraft/Joly files.]}
}
{
}

\Verse{170}
{
    \zR{到了潇湘馆内，一见黛玉灵柩，贾母 已哭得泪干气绝。}
}
{
    \eR{[English source not present in the checked-in Bencraft/Joly files.]}
}
{
}

\Verse{171}
{
    \zR{凤姐等再三劝住。}
}
{
    \eR{[English source not present in the checked-in Bencraft/Joly files.]}
}
{
}

\Verse{172}
{
    \zR{王夫人也哭了一场。}
}
{
    \eR{[English source not present in the checked-in Bencraft/Joly files.]}
}
{
}

\Verse{173}
{
    \zR{李纨便请贾母王夫人在里 间歇着，犹自落泪。}
}
{
    \eR{[English source not present in the checked-in Bencraft/Joly files.]}
}
{
}

\Verse{174}
{
    \zR{宝玉一到，想起未病之先，来到这里；今日屋在人亡，不禁嚎 啕大哭。}
}
{
    \eR{[English source not present in the checked-in Bencraft/Joly files.]}
}
{
}

\Verse{175}
{
    \zR{想起从前何等亲密，今日死别，怎不更加伤感!众人原恐宝玉病后过哀， 都来解劝。}
}
{
    \eR{[English source not present in the checked-in Bencraft/Joly files.]}
}
{
}

\Verse{176}
{
    \zR{宝玉已经哭得死去活来，大家搀扶歇息。}
}
{
    \eR{[English source not present in the checked-in Bencraft/Joly files.]}
}
{
}

\Verse{177}
{
    \zR{其馀随来的如宝钗，俱极痛哭。}
}
{
    \eR{[English source not present in the checked-in Bencraft/Joly files.]}
}
{
}

\Verse{178}
{
    \zR{独是宝玉必要叫紫鹃来见：“问明姑娘临死有何话说。”}
}
{
    \eR{[English source not present in the checked-in Bencraft/Joly files.]}
}
{
}

\Verse{179}
{
    \zR{紫鹃本来深恨宝玉，见如 此心里已回过来些，又有贾母王夫人都在这里，不敢洒落宝玉，便将林姑娘怎么复
        病，怎么烧毁帕子，焚化诗稿，并将临死说的话一一的都告诉了。}
}
{
    \eR{[English source not present in the checked-in Bencraft/Joly files.]}
}
{
}

\Verse{180}
{
    \zR{宝玉又哭得气噎 喉干。}
}
{
    \eR{[English source not present in the checked-in Bencraft/Joly files.]}
}
{
}

\Verse{181}
{
    \zR{探春趁便又将黛玉临终嘱咐带柩回南的话也说了一遍。}
}
{
    \eR{[English source not present in the checked-in Bencraft/Joly files.]}
}
{
}

\Verse{182}
{
    \zR{贾母王夫人又哭起 来。}
}
{
    \eR{[English source not present in the checked-in Bencraft/Joly files.]}
}
{
}

\Verse{183}
{
    \zR{多亏凤姐能言劝慰，略略止些，便请贾母等回去。}
}
{
    \eR{[English source not present in the checked-in Bencraft/Joly files.]}
}
{
}

\Verse{184}
{
    \zR{宝玉那里肯舍，无奈贾母逼 着，只得勉强回房。}
}
{
    \eR{[English source not present in the checked-in Bencraft/Joly files.]}
}
{
}

\Verse{185}
{
    \zR{贾母有了年纪的人，打从宝玉病起，日夜不宁，今又大痛一阵，已觉头晕身热， 虽是不放心惦着宝玉，却也扎挣不住，回到自己房中睡下。}
}
{
    \eR{[English source not present in the checked-in Bencraft/Joly files.]}
}
{
}

\Verse{186}
{
    \zR{王夫人更加心痛难禁， 也便回去，派了彩云帮着袭人照应，并说：“宝玉若再悲戚，速来告诉我们。”}
}
{
    \eR{[English source not present in the checked-in Bencraft/Joly files.]}
}
{
}

\Verse{187}
{
    \zR{宝 钗知是宝玉一时必不能舍，也不相劝，只用讽刺的话说他。}
}
{
    \eR{[English source not present in the checked-in Bencraft/Joly files.]}
}
{
}

\Verse{188}
{
    \zR{宝玉倒恐宝钗多心，也 便饮泣收心。}
}
{
    \eR{[English source not present in the checked-in Bencraft/Joly files.]}
}
{
}

\Verse{189}
{
    \zR{歇了一夜，倒也安稳。}
}
{
    \eR{[English source not present in the checked-in Bencraft/Joly files.]}
}
{
}

\Verse{190}
{
    \zR{明日一早，众人都来瞧他，但觉气虚身弱，心 病倒觉去了几分。}
}
{
    \eR{[English source not present in the checked-in Bencraft/Joly files.]}
}
{
}

\Verse{191}
{
    \zR{于是加意调养，渐渐的好起来。}
}
{
    \eR{[English source not present in the checked-in Bencraft/Joly files.]}
}
{
}

\Verse{192}
{
    \zR{贾母幸不成病，惟是王夫人心痛 未痊。}
}
{
    \eR{[English source not present in the checked-in Bencraft/Joly files.]}
}
{
}

\Verse{193}
{
    \zR{那日薛姨妈过来探望，看见宝玉精神略好，也就放心，暂且住下。}
}
{
    \eR{[English source not present in the checked-in Bencraft/Joly files.]}
}
{
}

\Verse{194}
{
    \zR{一日，贾母特请薛姨妈过去商量，说：“宝玉的命，都亏姨太太救的。}
}
{
    \eR{[English source not present in the checked-in Bencraft/Joly files.]}
}
{
}

\Verse{195}
{
    \zR{如今想 来不妨了。}
}
{
    \eR{[English source not present in the checked-in Bencraft/Joly files.]}
}
{
}

\Verse{196}
{
    \zR{独委屈了你的姑娘。}
}
{
    \eR{[English source not present in the checked-in Bencraft/Joly files.]}
}
{
}

\Verse{197}
{
    \zR{如今宝玉调养百日，身体复旧，又过了娘娘的功服， 正好圆房：要求姨太太作主，另择个上好的吉日。”}
}
{
    \eR{[English source not present in the checked-in Bencraft/Joly files.]}
}
{
}

\Verse{198}
{
    \zR{薛姨妈便道：“老太太主意很 好，何必问我？}
}
{
    \eR{[English source not present in the checked-in Bencraft/Joly files.]}
}
{
}

\Verse{199}
{
    \zR{宝丫头虽生的粗笨，心里却还是极明白的，他的情性老太太素日是 知道的。}
}
{
    \eR{[English source not present in the checked-in Bencraft/Joly files.]}
}
{
}

\Verse{200}
{
    \zR{但愿他们两口儿言和意顺，从此老太太也省好些心，我姐姐也安慰些，我 也放了心了。}
}
{
    \eR{[English source not present in the checked-in Bencraft/Joly files.]}
}
{
}

\Verse{201}
{
    \zR{老太太就定个日子。}
}
{
    \eR{[English source not present in the checked-in Bencraft/Joly files.]}
}
{
}

\Verse{202}
{
    \zR{还通知亲戚不用呢？”}
}
{
    \eR{[English source not present in the checked-in Bencraft/Joly files.]}
}
{
}

\Verse{203}
{
    \zR{贾母道：“宝玉和你们姑 娘生来第一件大事，况且费了多少周折，如今才得安逸，必要大家热闹几天。}
}
{
    \eR{[English source not present in the checked-in Bencraft/Joly files.]}
}
{
}

\Verse{204}
{
    \zR{亲戚 都要请的。}
}
{
    \eR{[English source not present in the checked-in Bencraft/Joly files.]}
}
{
}

\Verse{205}
{
    \zR{一来酬愿，二则咱们吃杯喜酒，也不枉我老人家操了好些心。”}
}
{
    \eR{[English source not present in the checked-in Bencraft/Joly files.]}
}
{
}

\Verse{206}
{
    \zR{薛姨妈 听着，自然也是喜欢的，便将要办妆奁的话也说了一番。}
}
{
    \eR{[English source not present in the checked-in Bencraft/Joly files.]}
}
{
}

\Verse{207}
{
    \zR{贾母道：“咱们亲上做亲， 我想也不必这么。}
}
{
    \eR{[English source not present in the checked-in Bencraft/Joly files.]}
}
{
}

\Verse{208}
{
    \zR{若说动用的，他屋里已经满了；必定宝丫头他心爱的要你几件， 姨太太就拿了来。}
}
{
    \eR{[English source not present in the checked-in Bencraft/Joly files.]}
}
{
}

\Verse{209}
{
    \zR{我看宝丫头也不是多心的人，比不的我那外孙女儿的脾气，所以 他不得长寿。”}
}
{
    \eR{[English source not present in the checked-in Bencraft/Joly files.]}
}
{
}

\Verse{210}
{
    \zR{说着，连薛姨妈也便落泪。}
}
{
    \eR{[English source not present in the checked-in Bencraft/Joly files.]}
}
{
}

\Verse{211}
{
    \zR{恰好凤姐进来，笑道：“老太太姑妈又 想着什么了？”}
}
{
    \eR{[English source not present in the checked-in Bencraft/Joly files.]}
}
{
}

\Verse{212}
{
    \zR{薛姨妈道：“我和老太太说起你林妹妹来，所以伤心。”}
}
{
    \eR{[English source not present in the checked-in Bencraft/Joly files.]}
}
{
}

\Verse{213}
{
    \zR{凤姐笑道： “老太太和姑妈且别伤心。}
}
{
    \eR{[English source not present in the checked-in Bencraft/Joly files.]}
}
{
}

\Verse{214}
{
    \zR{我刚才听了个笑话儿来了，意思说给老太太和姑妈听。”}
}
{
    \eR{[English source not present in the checked-in Bencraft/Joly files.]}
}
{
}

\Verse{215}
{
    \zR{贾母拭了拭眼泪，微笑道：“你又不知要编派谁呢？}
}
{
    \eR{[English source not present in the checked-in Bencraft/Joly files.]}
}
{
}

\Verse{216}
{
    \zR{你说来，我和姨太太听听。}
}
{
    \eR{[English source not present in the checked-in Bencraft/Joly files.]}
}
{
}

\Verse{217}
{
    \zR{说 不笑，我们可不依。”}
}
{
    \eR{[English source not present in the checked-in Bencraft/Joly files.]}
}
{
}

\Verse{218}
{
    \zR{只见那凤姐未从张口，先用两只手比着，笑弯了腰了。}
}
{
    \eR{[English source not present in the checked-in Bencraft/Joly files.]}
}
{
}

\Verse{219}
{
    \zR{未知他说出些什么来，下回分解。}
}
{
    \eR{[English source not present in the checked-in Bencraft/Joly files.]}
}
{
}

\Verse{220}
{
    \zR{(本章完)}
}
{
    \eR{[English source not present in the checked-in Bencraft/Joly files.]}
}
{
}
